% Claim proof file (REFACTOR TWIN) for row claim_3_19 -- "MarginalizingOutThe".
% This is the twin file produced during the `eqViaNodeMap_injective`
% refactor; the existing `claim_3_19_proof_MarginalizingOutThe.tex` stays
% untouched until cleanup at Phase 7 atomically renames this twin over it.
%
% Mathematically the LN-level proof is unchanged by the refactor: the
% lemma asserts the same CDMG-isomorphism
% $G_{\doit(W)} \cong \lp G_{\swig(W)} \rp^{\sm W^o}$ as before, with
% the same hypothesis $W \ins V$ and the same canonical bijection
% $\phi : J \cup V \to J \cup W^i \cup (V \sm W)$ defined by
% $v \mapsto v$ on $J \cup (V \sm W)$ and $w \mapsto w^i$ on $W$. The
% only addition versus the existing proof body is a new ``Injectivity
% of the carrier bijection on the source CDMG's node set'' paragraph
% inserted just before the closing ``Combining the four clauses
% (a)--(d) \ldots'' sentence, which underwrites the strengthened Lean
% predicate `refactor_eqViaNodeMap` introduced by refactor
% `eqViaNodeMap_injective` (which adds a 5th `Set.InjOn` conjunct on
% the source carrier $\uparrow G.J \cup \uparrow G.V$ to the original
% four image-equality conjuncts).
%
% Unlike claim_3_7's twin (where the new InjOn paragraph consumes the
% disjointness side condition $W_1 \cap W_2 = \emptyset$ to rule out
% two would-be cross-cell collisions), this row's InjOn paragraph
% needs no extra side condition beyond the lemma's standing hypothesis
% $W \ins V$: the source carrier $J \cup V$ partitions into three
% cells $\{J, V \sm W, W\}$, and pairwise disjointness of those cells
% follows from the CDMG axiom $J \cap V = \emptyset$ (def_3_1) and
% set algebra alone; across-cell collisions are then ruled out by the
% tagged-copy disjoint-union construction of def_3_12 (the SWIG
% definition), which realises $W^i$ as type-disjoint from the
% untagged image of $J \cup (V \sm W)$ in the SWIG carrier. The new
% paragraph lives at the LN/Lean interface and is not load-bearing
% for the LN-level mathematics of the lemma; the rest of the proof
% body -- statement-restatement block, structural observation on
% directed walks through $W^o$, and the four componentwise clauses
% (a)--(d) -- is verbatim from the existing
% `claim_3_19_proof_MarginalizingOutThe.tex`.
%
% This file is a sibling of `main.tex` inside `<subsection>/tex/`. The
% `subfiles` package lets it render both standalone and as part of
% `main.tex` via `\subfile{...}`.

\documentclass[main]{subfiles}

\begin{document}
% ref: claim_3_19 (proof, with statement restated for readability)
% title: MarginalizingOutThe
\def\rowref{claim\_3\_19}
\phantomsection\label{claim_3_19}

% --- statement (restated) ---------------------------------------------------
\begin{Lem}[Marginalizing out the output part of splitted nodes equals hard intervention]\label{marginalizing-out-the-output-part-of-splitted-nodes-equals-hard-intervention}
    Let $G = (J, V, E, L)$ be a CDMG (def \ref{def-cdmg}); throughout we use the notation of def \ref{not-cdmg} and recall the typing and axioms of def \ref{def-cdmg}, namely $J \cap V = \emptyset$, $E \ins (J \cup V) \times V$, $L \ins V \times V$, $L$ irreflexive ($(v_1, v_2) \in L \Rightarrow v_1 \neq v_2$), and $L$ symmetric ($(v_1, v_2) \in L \Leftrightarrow (v_2, v_1) \in L$). We rely throughout on the hard-intervention construction $\doit$ of def \ref{def:G_hard_intervention}, the node-splitting hard-intervention (SWIG) construction $\swig$ of def \ref{def:G_node-splitting_intervention}, and the marginalization operator $G^{\sm W}$ of def \ref{def:G_marginalization}. Let
    \[
        W \;\ins\; V
    \]
    be a subset of output nodes of $G$. The LN's phrasing ``subset of output nodes from $G$'' is read here as ``$W \ins V$'': by def \ref{def-cdmg} the set $V$ \emph{is} the output-node set of $G$, so the qualifier ``of output nodes'' is the standing typing $W \ins V$ and imposes no further constraint. The quantifiers ``for every CDMG $G$'' and ``for every $W \ins V$'' are read \emph{universally}; the degenerate case $W = \emptyset$ is admitted (then both sides of the displayed isomorphism collapse to $G$ verbatim, with the bijection $\phi$ below reducing to the identity on $J \cup V$, since $W = \emptyset$ makes $W^o = W^i = \emptyset$).

    \emph{Well-typedness of the two CDMGs in the displayed isomorphism.}
    \begin{itemize}
        \item \emph{LHS, $G_{\doit(W)}$.}
              The hard-intervention construction of def \ref{def:G_hard_intervention} is defined under the precondition $W \ins J \cup V$. Since $W \ins V$ and $V \ins J \cup V$, this precondition is met. By items~i.--iv.\ of def \ref{def:G_hard_intervention}, the four components of $G_{\doit(W)}$ are
              \[
                  J_{\doit(W)} \;=\; J \cup W,
                  \qquad
                  V_{\doit(W)} \;=\; V \sm W,
              \]
              \[
                  E_{\doit(W)} \;=\; E \sm \lC (v_1, v_2) \in E \st v_2 \in W \rC,
                  \qquad
                  L_{\doit(W)} \;=\; L \sm \lC (v_1, v_2) \in L \st v_2 \in W \rC,
              \]
              read along the symmetrised reading of item~iv.\ of def \ref{def:G_hard_intervention}. The carrier of $G_{\doit(W)}$ is
              \[
                  J_{\doit(W)} \cup V_{\doit(W)} \;=\; (J \cup W) \cup (V \sm W) \;=\; J \cup V,
              \]
              using $W \ins V$ and $J \cap V = \emptyset$.

        \item \emph{RHS, $\lp G_{\swig(W)} \rp^{\sm W^o}$.}
              The SWIG construction of def \ref{def:G_node-splitting_intervention} requires the input CDMG $G$ to be a CADMG (i.e.\ acyclic in the sense of def \ref{def-acylic}) and $W \ins V$. The LN's literal hypothesis on $G$ in the present row, however, asserts only that $G$ is a CDMG; we therefore explicitly require here, in addition to the standing CDMG hypothesis on $G$, the CDMG-level reading of $\swig(W)$ as the set-builder construction of items~i.--iv.\ of def \ref{def:G_node-splitting_intervention} (the construction is well-typed at the CDMG level using only the CDMG axioms, per the ``Paradigm observation (which axioms each clause uses)'' paragraph of def \ref{def:G_node-splitting_intervention}; acyclicity of the input is not consumed by items~i.--iv.\ themselves, only by the downstream upgrade to a CADMG, which is not needed here since both sides of the displayed isomorphism are read at the CDMG level). With $W \ins V$, items~i.--iv.\ of def \ref{def:G_node-splitting_intervention} yield
              \[
                  J_{\swig(W)} \;=\; J \dcup W^i,
                  \qquad
                  V_{\swig(W)} \;=\; (V \sm W) \dcup W^o,
              \]
              \[
                  E_{\swig(W)} \;=\; \lC (v_1^i, v_2^o) \st (v_1, v_2) \in E \rC,
                  \qquad
                  L_{\swig(W)} \;=\; \lC (v_1^o, v_2^o) \st (v_1, v_2) \in L \rC,
              \]
              with $W^o := \lC w^o \st w \in W \rC$ and $W^i := \lC w^i \st w \in W \rC$ the two type-level tagged copies of $W$ from def \ref{def:G_node-splitting_intervention} (so $W^o \cap W^i = \emptyset$, and both $W^o, W^i$ are disjoint from every element of $J \cup V$), and with the unsplit-injection shorthand ``$v^o := v^i := v$ for $v \in J \cup (V \sm W)$'' from def \ref{def:G_node-splitting_intervention} in force inside the set-builders for $E_{\swig(W)}$ and $L_{\swig(W)}$.

              For the outer marginalization $(\,\cdot\,)^{\sm W^o}$ applied to $G_{\swig(W)}$ to be defined per def \ref{def:G_marginalization}'s precondition ``$W' \ins V'$ of the input CDMG'' (instantiated with $W' = W^o$ and $V' = V_{\swig(W)} = (V \sm W) \dcup W^o$), we need $W^o \ins V_{\swig(W)}$. This is immediate from $W^o \ins (V \sm W) \dcup W^o = V_{\swig(W)}$. By items~i.\ and~ii.\ of def \ref{def:G_marginalization} applied to the input CDMG $G_{\swig(W)}$ with $W^o$ marginalized out, the input-node and output-node sets of $\lp G_{\swig(W)} \rp^{\sm W^o}$ are
              \[
                  J_{\lp G_{\swig(W)} \rp^{\sm W^o}} \;=\; J_{\swig(W)} \;=\; J \dcup W^i,
                  \qquad
                  V_{\lp G_{\swig(W)} \rp^{\sm W^o}} \;=\; V_{\swig(W)} \sm W^o \;=\; V \sm W,
              \]
              where the second equality uses $(V \sm W) \cap W^o = \emptyset$ (since $W^o$ is type-level disjoint from $V$ by the tagged-copy construction). The carrier of $\lp G_{\swig(W)} \rp^{\sm W^o}$ is therefore
              \[
                  J_{\lp G_{\swig(W)} \rp^{\sm W^o}} \cup V_{\lp G_{\swig(W)} \rp^{\sm W^o}} \;=\; (J \dcup W^i) \cup (V \sm W) \;=\; J \cup W^i \cup (V \sm W).
              \]
              The directed-edge and bidirected-edge sets $E_{\lp G_{\swig(W)} \rp^{\sm W^o}}$ and $L_{\lp G_{\swig(W)} \rp^{\sm W^o}}$ are obtained from items~iii.\ and~iv.\ of def \ref{def:G_marginalization} applied to $G_{\swig(W)}$ marginalized at $W^o$; their explicit unfolding into directed-walk-through-$W^o$ and bifurcation-through-$W^o$ predicates ($\Phi_E$ and $\Phi_L$ of def \ref{def:G_marginalization}) is deferred to the proof.
    \end{itemize}

    \emph{The conclusion: a CDMG-isomorphism with explicit bijection $\phi$.}
    There exists a CDMG-isomorphism
    \[
        \phi \;:\; G_{\doit(W)} \;\xrightarrow{\;\cong\;}\; \lp G_{\swig(W)} \rp^{\sm W^o}
    \]
    in the sense made precise below, defined on the carrier $J_{\doit(W)} \cup V_{\doit(W)} = J \cup V$ of $G_{\doit(W)}$ by the following case split (exhaustive and pairwise disjoint because $W \ins V$ and $J \cap V = \emptyset$, hence $J$, $V \sm W$, $W$ partition $J \cup V$):
    \[
        \phi(v) \;:=\;
        \begin{cases}
            v & \text{if } v \in J, \\
            v & \text{if } v \in V \sm W, \\
            v^i & \text{if } v \in W.
        \end{cases}
    \]
    Equivalently, $\phi$ is the identity on $J \cup (V \sm W)$ and sends $w \mapsto w^i$ on $W$ (the LN's literal ``$w \mapsto w^i$'' clause, supplemented with the identity action on $J \cup (V \sm W)$ that the LN leaves implicit). The codomain of $\phi$ is the carrier $J_{\lp G_{\swig(W)} \rp^{\sm W^o}} \cup V_{\lp G_{\swig(W)} \rp^{\sm W^o}} = J \cup W^i \cup (V \sm W)$ of $\lp G_{\swig(W)} \rp^{\sm W^o}$, displayed above; on each branch of the case split, $\phi(v)$ lies in this codomain (for $v \in J$: $\phi(v) = v \in J$; for $v \in V \sm W$: $\phi(v) = v \in V \sm W$; for $v \in W$: $\phi(v) = v^i \in W^i$). Bijectivity of $\phi$ holds because the codomain partitions as $J \cup W^i \cup (V \sm W)$ with $\phi$ restricting to a bijection on each of the three partition pieces (the identity on $J$, the identity on $V \sm W$, and the injection $w \mapsto w^i$ from $W$ to $W^i$, which is a bijection by def \ref{def:G_node-splitting_intervention}'s tagged-copy construction).

    \emph{Reading of ``CDMG-isomorphism''.}
    The relation $G_{\doit(W)} \cong \lp G_{\swig(W)} \rp^{\sm W^o}$ asserted by the LN's $\cong$ symbol is read here as the conjunction of the following four clauses, jointly stating that $\phi$ above is a CDMG-isomorphism of the two CDMGs (def \ref{def-cdmg}) under their four-tuple structures $(J, V, E, L)$:
    \begin{enumerate}[label=(\alph*)]
        \item \emph{$\phi$ maps input nodes to input nodes.}
              For every $v \in J_{\doit(W)} = J \cup W$,
              \[
                  \phi(v) \;\in\; J_{\lp G_{\swig(W)} \rp^{\sm W^o}} \;=\; J \dcup W^i.
              \]
              Equivalently, the restriction $\phi\rvert_{J \cup W} : J \cup W \to J \dcup W^i$ is a bijection: on $J$ it is the identity (and $J \ins J \dcup W^i$), and on $W$ it sends $w \mapsto w^i \in W^i$ (and the latter is a bijection $W \to W^i$ by def \ref{def:G_node-splitting_intervention}).

        \item \emph{$\phi$ maps output nodes to output nodes.}
              For every $v \in V_{\doit(W)} = V \sm W$,
              \[
                  \phi(v) \;\in\; V_{\lp G_{\swig(W)} \rp^{\sm W^o}} \;=\; V \sm W.
              \]
              Equivalently, the restriction $\phi\rvert_{V \sm W} : V \sm W \to V \sm W$ is the identity, hence trivially a bijection.

        \item \emph{$\phi$ preserves directed edges.}
              For every $(v_1, v_2) \in (J \cup V) \times V$ that lies in the typing of $E_{\doit(W)}$ (i.e.\ $v_1 \in J_{\doit(W)} \cup V_{\doit(W)} = J \cup V$ and $v_2 \in V_{\doit(W)} = V \sm W$, matching the precondition $E_{\doit(W)} \ins (J_{\doit(W)} \cup V_{\doit(W)}) \times V_{\doit(W)}$ of def \ref{def-cdmg} applied to $G_{\doit(W)}$),
              \[
                  (v_1, v_2) \;\in\; E_{\doit(W)}
                  \quad\Longleftrightarrow\quad
                  (\phi(v_1), \phi(v_2)) \;\in\; E_{\lp G_{\swig(W)} \rp^{\sm W^o}}.
              \]

        \item \emph{$\phi$ preserves bidirected edges.}
              For every $(v_1, v_2) \in V \times V$ that lies in the typing of $L_{\doit(W)}$ (i.e.\ $v_1, v_2 \in V_{\doit(W)} = V \sm W$, matching the precondition $L_{\doit(W)} \ins V_{\doit(W)} \times V_{\doit(W)}$ of def \ref{def-cdmg} applied to $G_{\doit(W)}$),
              \[
                  (v_1, v_2) \;\in\; L_{\doit(W)}
                  \quad\Longleftrightarrow\quad
                  (\phi(v_1), \phi(v_2)) \;\in\; L_{\lp G_{\swig(W)} \rp^{\sm W^o}}.
              \]
    \end{enumerate}
    Clauses~(a) and~(b) together state that $\phi$ is a bijection of carriers respecting the input/output partition (so it transports the $J$- and $V$-components of the two CDMGs onto each other); clauses~(c) and~(d) state that the directed- and bidirected-edge sets of the two CDMGs correspond under the bijection $\phi$, where the edge sets on the right-hand side are the marginalization-of-SWIG edge sets $E_{\lp G_{\swig(W)} \rp^{\sm W^o}}$ and $L_{\lp G_{\swig(W)} \rp^{\sm W^o}}$ of def \ref{def:G_marginalization} applied to $G_{\swig(W)}$ with $W^o$ marginalized out (unfolded into the directed-walk-through-$W^o$ predicate $\Phi_E$ and the bifurcation-through-$W^o$ predicate $\Phi_L$ at items~iii.\ and~iv.\ of def \ref{def:G_marginalization}). The conjunctive unpacking of clauses~(a)--(d) into the four field-by-field bijection / preservation statements of $J$, $V$, $E$, and $L$ is deferred to the proof.
\end{Lem}

% --- proof ------------------------------------------------------------------
\begin{proof}
% Adapted from the LN's own \Claude{}-wrapped proof block at
% lecture-notes/lecture_notes/graphs.tex:1176-1211, with four structural
% edits:
%   - the LN proof writes the isomorphism in the inverse direction
%     ($w^i \mapsto w$); we run the forward direction
%     ($w \mapsto w^i$) to match the canonical statement's $\phi$;
%   - the LN's three sections ``Node sets / Directed edges / Bidirected
%     edges'' are repackaged as the canonical statement's four clauses
%     (a)--(d) of CDMG-isomorphism (input bijection, output bijection,
%     directed-edge biconditional, bidirected-edge biconditional);
%   - the LN's two implicit shortcuts -- ``no directed walk can pass
%     through $W^o$'' and ``bifurcations with intermediates in $W^o$
%     collapse'' -- are spelled out inline against the explicit edge-set
%     of $E_{\swig(W)}$ from def_3_12, item iii.;
%   - a new ``fifth componentwise check'' paragraph -- injectivity of
%     the canonical carrier bijection $\phi$ on the source CDMG's joint
%     input-output set $J_{\doit(W)} \cup V_{\doit(W)} = J \cup V$ --
%     is inserted at the end of the componentwise verification, just
%     before the closing ``Combining the four clauses (a)--(d) \ldots''
%     sentence.  The new paragraph is \emph{not} in the LN's \Claude
%     proof (the LN identifies the carriers implicitly via the
%     case-split definition of $\phi$ and the bijectivity-on-pieces
%     remark in the statement block, so no source-side injectivity
%     certificate is needed in the LN's reading); it is added here to
%     underwrite the strengthened Lean-encoding predicate
%     `refactor_eqViaNodeMap` introduced by refactor
%     `eqViaNodeMap_injective` (whose 5th conjunct is exactly a
%     `Set.InjOn` on the source carrier $\uparrow G.J \cup \uparrow
%     G.V$).  Mathematically it is a straightforward three-cell case
%     analysis on the source $J \cup V$ using only the standing CDMG
%     axiom $J \cap V = \emptyset$ (def_3_1) and def_3_12's tagged-copy
%     disjoint-union construction -- no extra side condition (such as
%     a disjointness hypothesis between two split sets) is consumed,
%     in contrast to claim_3_7's twin where the new InjOn paragraph
%     consumes the disjointness side condition $W_1 \cap W_2 =
%     \emptyset$.  The new paragraph lives at the LN/Lean interface
%     and is not load-bearing for the LN-level mathematics of the
%     lemma.
We establish the four clauses (a)--(d) of CDMG-isomorphism for the bijection $\phi : J \cup V \to J \cup W^i \cup (V \sm W)$ of the statement, defined by $\phi(v) = v$ for $v \in J \cup (V \sm W)$ and $\phi(w) = w^i$ for $w \in W$ (def \ref{def:G_node-splitting_intervention}'s tagged-copy construction supplies the $W \to W^i$ bijection $w \mapsto w^i$). Clauses (a) and (b) are immediate from the case-split definition of $\phi$ and the carrier identifications already worked out in the statement; clauses (c) and (d) require unfolding the marginalization-through-$W^o$ predicates $\Phi_E$ and $\Phi_L$ of def \ref{def:G_marginalization} applied to $G_{\swig(W)}$, and crucially rely on the structural observation that no directed walk in $G_{\swig(W)}$ can pass through $W^o$ as an intermediate node. We prove that observation first, then work the four clauses in turn.

\medskip\noindent\emph{Structural observation: no directed walk in $G_{\swig(W)}$ passes through $W^o$ as an intermediate node.}
By def \ref{def:G_node-splitting_intervention}, item~iii.,
\[
    E_{\swig(W)} \;=\; \lC (v_1^i, v_2^o) \st (v_1, v_2) \in E \rC,
\]
so every directed edge of $G_{\swig(W)}$ has source of the form $v_1^i$ (lying in $W^i$ if $v_1 \in W$, untagged in $J \cup (V \sm W)$ otherwise, by the shorthand $v^i := v$ of def \ref{def:G_node-splitting_intervention} for $v \in J \cup (V \sm W)$) and target of the form $v_2^o$ (lying in $W^o$ if $v_2 \in W$, untagged in $V \sm W$ otherwise, by the shorthand $v^o := v$ for $v \in V \sm W$). Consequently \emph{no directed edge of $G_{\swig(W)}$ has its source in $W^o$}: a putative such edge $(w^o, y) \in E_{\swig(W)}$ would force, by the set-builder, $w^o = v_1^i$ for some $v_1 \in J \cup V$, but $w^o \in W^o$ while $v_1^i \in W^i \cup J \cup (V \sm W)$, two type-disjoint pieces (the tagged-copy construction of def \ref{def:G_node-splitting_intervention} makes $W^o \cap W^i = \emptyset$ and $W^o \cap (J \cup (V \sm W)) = \emptyset$).

Now suppose $\pi = (b_0, b_1, \ldots, b_m)$ is a directed walk in $G_{\swig(W)}$ (in the sense of def \ref{def:G_marginalization}, item~iii., clauses~(a)--(c) applied to $G_{\swig(W)}$, namely $m \ge 1$ and $(b_{i-1}, b_i) \in E_{\swig(W)}$ for every $i \in \{1, \ldots, m\}$) and that some intermediate node $b_j$ ($1 \le j \le m-1$) lies in $W^o$. Then the next edge $(b_j, b_{j+1}) \in E_{\swig(W)}$ has source $b_j \in W^o$ -- but by the previous paragraph no directed edge of $G_{\swig(W)}$ has its source in $W^o$. Contradiction. Hence no directed walk in $G_{\swig(W)}$ has any intermediate node in $W^o$.

Corollary: \emph{the only directed walks in $G_{\swig(W)}$ whose intermediates lie in $W^o$ are the single-edge walks ($m = 1$).} The clause ``intermediates in $W^o$'' is vacuous when $m = 1$ (the index range $\{1, \ldots, m-1\}$ is empty), so any single-edge walk $(b_0, b_1)$ with $(b_0, b_1) \in E_{\swig(W)}$ qualifies; longer walks ($m \ge 2$) have at least one intermediate $b_1, \ldots, b_{m-1}$ which, by the contradiction above, cannot lie in $W^o$.

\medskip\noindent\textbf{Clause (a) -- $\phi$ maps input nodes to input nodes.}
For $v \in J_{\doit(W)} = J \cup W$: if $v \in J$ then $\phi(v) = v \in J \ins J \dcup W^i = J_{\lp G_{\swig(W)} \rp^{\sm W^o}}$; if $v \in W$ then $\phi(v) = v^i \in W^i \ins J \dcup W^i$. In either case $\phi(v) \in J_{\lp G_{\swig(W)} \rp^{\sm W^o}}$. The restriction $\phi\rvert_{J \cup W} : J \cup W \to J \dcup W^i$ is a bijection because $J \cup W$ partitions as $J \dcup W$ (since $J \cap V = \emptyset$ and $W \ins V$ give $J \cap W = \emptyset$), $J \dcup W^i$ partitions as $J \dcup W^i$ by construction (def \ref{def:G_node-splitting_intervention}: $W^i$ is type-disjoint from $J$), and $\phi$ restricts to the identity bijection $J \to J$ on $J$ and to the bijection $W \to W^i$ $(w \mapsto w^i)$ on $W$ (the latter by def \ref{def:G_node-splitting_intervention}'s tagged-copy construction). \checkmark

\medskip\noindent\textbf{Clause (b) -- $\phi$ maps output nodes to output nodes.}
For $v \in V_{\doit(W)} = V \sm W$: $\phi(v) = v$ (by the case split, since $V \sm W \ins V \sm W$), and $v \in V \sm W = V_{\lp G_{\swig(W)} \rp^{\sm W^o}}$. The restriction $\phi\rvert_{V \sm W} : V \sm W \to V \sm W$ is the identity, hence trivially a bijection. \checkmark

\medskip\noindent\textbf{Clause (c) -- $\phi$ preserves directed edges.}
Let $(v_1, v_2) \in (J \cup V) \times (V \sm W)$ be a pair lying in the typing of $E_{\doit(W)}$ (so $v_1 \in J \cup V$ and $v_2 \in V \sm W$, matching the precondition of def \ref{def-cdmg} for $G_{\doit(W)}$). Under $\phi$, the image pair $(\phi(v_1), \phi(v_2))$ has:
\begin{itemize}
    \item $\phi(v_2) = v_2$, since $v_2 \in V \sm W$;
    \item $\phi(v_1) = v_1$ if $v_1 \in J \cup (V \sm W)$, and $\phi(v_1) = v_1^i \in W^i$ if $v_1 \in W$.
\end{itemize}

We compare the two biconditional clauses.

\smallskip\noindent\emph{LHS: $(v_1, v_2) \in E_{\doit(W)}$.} By def \ref{def:G_hard_intervention}, item~iii.,
\[
    E_{\doit(W)} \;=\; E \;\sm\; \lC (v_1', v_2') \in E \st v_2' \in W \rC,
\]
so $(v_1, v_2) \in E_{\doit(W)}$ iff $(v_1, v_2) \in E$ and $v_2 \notin W$. The second condition $v_2 \notin W$ is automatic from the typing $v_2 \in V \sm W$. Hence
\[
    (v_1, v_2) \in E_{\doit(W)} \;\Longleftrightarrow\; (v_1, v_2) \in E.
\]

\smallskip\noindent\emph{RHS: $(\phi(v_1), \phi(v_2)) \in E_{\lp G_{\swig(W)} \rp^{\sm W^o}}$.} By def \ref{def:G_marginalization}, item~iii., applied to $G_{\swig(W)}$ with $W^o$ marginalized out,
\[
    E_{\lp G_{\swig(W)} \rp^{\sm W^o}} \;=\; \lC (\ul{v}, \ol{v}) \in \lp J_{\swig(W)} \cup (V_{\swig(W)} \sm W^o) \rp \times (V_{\swig(W)} \sm W^o) \st \Phi_E^{\swig(W), W^o}(\ul{v}, \ol{v}) \rC,
\]
where $\Phi_E^{\swig(W), W^o}(\ul{v}, \ol{v})$ asserts the existence of a directed walk in $G_{\swig(W)}$ from $\ul{v}$ to $\ol{v}$ of length $\ge 1$ with intermediate nodes in $W^o$. By the structural observation above, every such walk is a single-edge walk: there exist no walks in $G_{\swig(W)}$ of length $\ge 2$ with intermediate nodes in $W^o$. Hence
\[
    \Phi_E^{\swig(W), W^o}(\ul{v}, \ol{v}) \;\Longleftrightarrow\; (\ul{v}, \ol{v}) \in E_{\swig(W)},
\]
and so
\[
    E_{\lp G_{\swig(W)} \rp^{\sm W^o}} \;=\; \lC (\ul{v}, \ol{v}) \in \lp J_{\swig(W)} \cup (V_{\swig(W)} \sm W^o) \rp \times (V_{\swig(W)} \sm W^o) \st (\ul{v}, \ol{v}) \in E_{\swig(W)} \rC.
\]

Now $(\phi(v_1), \phi(v_2)) \in E_{\lp G_{\swig(W)} \rp^{\sm W^o}}$ holds iff $(\phi(v_1), \phi(v_2)) \in E_{\swig(W)}$, the typing constraints being met automatically: $\phi(v_2) = v_2 \in V \sm W = V_{\swig(W)} \sm W^o$ (output-side restriction satisfied), and $\phi(v_1)$ lies in $J \cup W^i \cup (V \sm W) = J_{\swig(W)} \cup (V_{\swig(W)} \sm W^o)$ by the case split. By def \ref{def:G_node-splitting_intervention}, item~iii.,
\[
    E_{\swig(W)} \;=\; \lC (u_1^i, u_2^o) \st (u_1, u_2) \in E \rC,
\]
so $(\phi(v_1), \phi(v_2)) \in E_{\swig(W)}$ iff there exist $(u_1, u_2) \in E$ with $\phi(v_1) = u_1^i$ and $\phi(v_2) = u_2^o$. We show this is equivalent to $(v_1, v_2) \in E$.

\smallskip\noindent($\Leftarrow$.) Suppose $(v_1, v_2) \in E$. Set $u_1 := v_1$ and $u_2 := v_2$. Then $(u_1, u_2) \in E$ and:
\begin{itemize}
    \item If $v_1 \in J \cup (V \sm W)$: $u_1^i = v_1^i = v_1 = \phi(v_1)$ by the shorthand $v^i := v$ of def \ref{def:G_node-splitting_intervention}.
    \item If $v_1 \in W$: $u_1^i = v_1^i = \phi(v_1)$ directly.
    \item For $v_2 \in V \sm W$: $u_2^o = v_2^o = v_2 = \phi(v_2)$ by the shorthand $v^o := v$ for $v \in V \sm W$.
\end{itemize}
So $(\phi(v_1), \phi(v_2)) = (u_1^i, u_2^o) \in E_{\swig(W)}$.

\smallskip\noindent($\Rightarrow$.) Suppose $(\phi(v_1), \phi(v_2)) = (u_1^i, u_2^o)$ for some $(u_1, u_2) \in E$. We must show $(v_1, v_2) \in E$. Recover $v_2$ from $u_2$: $\phi(v_2) = v_2$ (since $v_2 \in V \sm W$) and $\phi(v_2) = u_2^o$, so $v_2 = u_2^o$. Since $v_2 \in V \sm W$ (untagged), the equality $v_2 = u_2^o$ forces $u_2 \in V \sm W$ (else $u_2 \in W$ would give $u_2^o \in W^o$, type-disjoint from $V \sm W$, contradicting $v_2 \in V \sm W$); under that branch the shorthand gives $u_2^o = u_2$, so $u_2 = v_2$. Recover $v_1$ from $u_1$: $\phi(v_1) = u_1^i$. Two cases on the case split for $\phi(v_1)$:
\begin{itemize}
    \item If $v_1 \in J \cup (V \sm W)$, then $\phi(v_1) = v_1$ (untagged in $J \cup (V \sm W)$); the equation $v_1 = u_1^i$ then forces $u_1 \in J \cup (V \sm W)$ (else $u_1 \in W$ gives $u_1^i \in W^i$, type-disjoint from $J \cup (V \sm W)$); under that branch $u_1^i = u_1$, so $u_1 = v_1$.
    \item If $v_1 \in W$, then $\phi(v_1) = v_1^i \in W^i$; the equation $v_1^i = u_1^i$ forces $u_1 \in W$ (since $v \mapsto v^i$ restricted to $W$ is the bijection $W \to W^i$ of def \ref{def:G_node-splitting_intervention} and $W^i$ is type-disjoint from the untagged shorthand image $J \cup (V \sm W)$); under that branch $u_1^i = u_1^i$ and the injection $W \to W^i$ gives $u_1 = v_1$.
\end{itemize}
In both sub-cases $(u_1, u_2) = (v_1, v_2)$, so $(v_1, v_2) = (u_1, u_2) \in E$.

\smallskip\noindent\emph{Conclusion of clause (c).} Combining the LHS and RHS unfoldings,
\[
    (v_1, v_2) \in E_{\doit(W)}
        \;\Leftrightarrow\; (v_1, v_2) \in E
        \;\Leftrightarrow\; (\phi(v_1), \phi(v_2)) \in E_{\swig(W)}
        \;\Leftrightarrow\; (\phi(v_1), \phi(v_2)) \in E_{\lp G_{\swig(W)} \rp^{\sm W^o}}.
\]
\checkmark

\medskip\noindent\textbf{Clause (d) -- $\phi$ preserves bidirected edges.}
Let $(v_1, v_2) \in (V \sm W) \times (V \sm W)$ lie in the typing of $L_{\doit(W)}$ (so $v_1, v_2 \in V \sm W$; def \ref{def-cdmg}'s irreflexivity of $L_{\doit(W)}$ further forces $v_1 \neq v_2$ when there is an edge between them, but we keep the typing as $(V \sm W) \times (V \sm W)$ and let the biconditional handle the irreflexivity automatically). Under $\phi$ the image pair is $(\phi(v_1), \phi(v_2)) = (v_1, v_2)$, since both $v_1, v_2 \in V \sm W$ lie in the identity branch of the case split.

\smallskip\noindent\emph{LHS: $(v_1, v_2) \in L_{\doit(W)}$.} By def \ref{def:G_hard_intervention}, item~iv., read along the symmetrised reading flagged in the remark following item~iv.\ (cf.\ the remark on item~iv.\ and the symmetry of $L$ in def \ref{def:G_hard_intervention} and the registered deviation flagged in \refrow{claim_3_4}'s proof),
\[
    L_{\doit(W)} \;=\; \lC (v_1', v_2') \in L \st v_1' \notin W \,\wedge\, v_2' \notin W \rC.
\]
For our $(v_1, v_2)$ with $v_1, v_2 \in V \sm W$, both $v_1 \notin W$ and $v_2 \notin W$ are automatic from the typing, so
\[
    (v_1, v_2) \in L_{\doit(W)} \;\Longleftrightarrow\; (v_1, v_2) \in L.
\]

\smallskip\noindent\emph{RHS: $(v_1, v_2) \in L_{\lp G_{\swig(W)} \rp^{\sm W^o}}$.} By def \ref{def:G_marginalization}, item~iv., applied to $G_{\swig(W)}$ with $W^o$ marginalized out,
\[
    L_{\lp G_{\swig(W)} \rp^{\sm W^o}} \;=\; \lC (\ul{v}, \ol{v}) \in (V_{\swig(W)} \sm W^o) \times (V_{\swig(W)} \sm W^o) \st \ul{v} \neq \ol{v} \,\wedge\, \Phi_L^{\swig(W), W^o}(\ul{v}, \ol{v}) \rC,
\]
where $\Phi_L^{\swig(W), W^o}(\ul{v}, \ol{v})$ asserts the existence of a bifurcation in $G_{\swig(W)}$ between $\ul{v}$ and $\ol{v}$ with intermediate nodes in $W^o$, in the sense of def \ref{def:G_marginalization}, item~iv. We unfold $\Phi_L^{\swig(W), W^o}$ using the structural observation above and show it is equivalent to $(\ul{v}, \ol{v}) \in L_{\swig(W)}$.

A bifurcation in $G_{\swig(W)}$ between $\ul{v}$ and $\ol{v}$ with intermediates in $W^o$ is, per def \ref{def:G_marginalization}, item~iv., a tuple $(w_0, \ldots, w_n)$ with hinge index $k \in \{1, \ldots, n\}$ such that $w_0 = \ul{v}$, $w_n = \ol{v}$, all intermediates $w_1, \ldots, w_{n-1} \in W^o$ (vacuous when $n = 1$), and:
\begin{itemize}
    \item Left directed arm (vacuous when $k = 1$): $(w_i, w_{i-1}) \in E_{\swig(W)}$ for every $i \in \{1, \ldots, k-1\}$.
    \item Central hinge at index $k$: $(w_k, w_{k-1}) \in E_{\swig(W)}$ or $(w_{k-1}, w_k) \in L_{\swig(W)}$ (with the directed alternative excluded when $k = n$, by clause~(f) of def \ref{def:G_marginalization}, item~iv.).
    \item Right directed arm (vacuous when $k = n$): $(w_i, w_{i+1}) \in E_{\swig(W)}$ for every $i \in \{k, \ldots, n-1\}$.
\end{itemize}

\emph{Claim: every bifurcation in $G_{\swig(W)}$ with intermediates in $W^o$ collapses to a single bidirected hinge edge, i.e.\ $n = 1$, $k = 1$, and the witness is the bidirected hinge alternative $(w_0, w_1) \in L_{\swig(W)}$ with $w_0 = \ul{v}$, $w_1 = \ol{v}$.}

\noindent Proof of claim. Suppose for contradiction $n \ge 2$ (so there is at least one intermediate index $j \in \{1, \ldots, n-1\}$ with $w_j \in W^o$). Fix such a $j$. We exhibit a directed edge of $G_{\swig(W)}$ with source $w_j$, contradicting the structural observation. Three cases for the position of $j$ relative to the hinge index $k$:
\begin{itemize}
    \item \emph{Left of the hinge ($j \le k - 1$, i.e.\ $j \in \{1, \ldots, k-1\}$).} The left-arm clause at index $i = j$ gives $(w_j, w_{j-1}) \in E_{\swig(W)}$ -- a directed edge with source $w_j \in W^o$.
    \item \emph{At the hinge ($j = k$).} Since $j \le n - 1$, we have $k = j \le n - 1$, so the right-arm clause at index $i = k$ is non-vacuous and gives $(w_k, w_{k+1}) \in E_{\swig(W)}$ -- a directed edge with source $w_k = w_j \in W^o$.
    \item \emph{Right of the hinge ($j \ge k + 1$, i.e.\ $j \in \{k+1, \ldots, n-1\}$).} The right-arm clause at index $i = j$ gives $(w_j, w_{j+1}) \in E_{\swig(W)}$ -- a directed edge with source $w_j \in W^o$.
\end{itemize}
(The three cases exhaust $j \in \{1, \ldots, n-1\}$ since $k \in \{1, \ldots, n\}$.) In every case the bifurcation skeleton forces a directed edge of $G_{\swig(W)}$ with source $w_j \in W^o$, which the structural observation excludes. Hence no $n \ge 2$ bifurcation in $G_{\swig(W)}$ with intermediates in $W^o$ exists.

So $n = 1$. With $n = 1$, the intermediate-index range $\{1, \ldots, n-1\} = \emptyset$ is vacuous (no intermediates), and the only valid hinge index is $k \in \{1, \ldots, n\} = \{1\}$, so $k = 1 = n$. By clause~(f) of def \ref{def:G_marginalization}, item~iv., when $k = n$ only the bidirected hinge alternative is admitted: $(w_0, w_1) \in L_{\swig(W)}$ (the directed alternative $(w_1, w_0) \in E_{\swig(W)}$ is excluded at $k = n$, since it would place the only arrowhead at $w_n$ at the wrong node). With $w_0 = \ul{v}$ and $w_1 = \ol{v}$, this is the single bidirected edge $(\ul{v}, \ol{v}) \in L_{\swig(W)}$. \emph{Claim proven.}

By the claim,
\[
    \Phi_L^{\swig(W), W^o}(\ul{v}, \ol{v}) \;\Longleftrightarrow\; (\ul{v}, \ol{v}) \in L_{\swig(W)},
\]
and therefore
\[
    L_{\lp G_{\swig(W)} \rp^{\sm W^o}} \;=\; \lC (\ul{v}, \ol{v}) \in (V_{\swig(W)} \sm W^o) \times (V_{\swig(W)} \sm W^o) \st \ul{v} \neq \ol{v} \,\wedge\, (\ul{v}, \ol{v}) \in L_{\swig(W)} \rC.
\]

Now apply this with $(\ul{v}, \ol{v}) = (\phi(v_1), \phi(v_2)) = (v_1, v_2)$. The typing $\ul{v}, \ol{v} \in V_{\swig(W)} \sm W^o = V \sm W$ is met by hypothesis. The condition $\ul{v} \neq \ol{v}$ is the irreflexivity constraint of def \ref{def:G_marginalization}, item~iv.; conversely it is encoded in $L$ via the irreflexivity of $L$ (def \ref{def-cdmg}), so $(v_1, v_2) \in L \Rightarrow v_1 \neq v_2$ automatically, and the constraint $v_1 \neq v_2$ on the RHS is matched by the LHS's $(v_1, v_2) \in L$'s irreflexivity (an edge in $L$ forces $v_1 \neq v_2$). Thus
\[
    (v_1, v_2) \in L_{\lp G_{\swig(W)} \rp^{\sm W^o}} \;\Longleftrightarrow\; v_1 \neq v_2 \,\wedge\, (v_1, v_2) \in L_{\swig(W)}.
\]

It remains to show $(v_1, v_2) \in L_{\swig(W)} \Leftrightarrow (v_1, v_2) \in L$, for $v_1, v_2 \in V \sm W$. By def \ref{def:G_node-splitting_intervention}, item~iv.,
\[
    L_{\swig(W)} \;=\; \lC (u_1^o, u_2^o) \st (u_1, u_2) \in L \rC.
\]
\smallskip\noindent($\Leftarrow$.) If $(v_1, v_2) \in L$, set $u_1 := v_1, u_2 := v_2$; then $(u_1, u_2) \in L$ and $u_k^o = v_k^o = v_k$ for $k = 1, 2$ by the shorthand $v^o := v$ on $V \sm W$. So $(v_1, v_2) = (u_1^o, u_2^o) \in L_{\swig(W)}$.

\smallskip\noindent($\Rightarrow$.) Conversely, suppose $(v_1, v_2) = (u_1^o, u_2^o)$ for some $(u_1, u_2) \in L$. The endpoints $v_1, v_2 \in V \sm W$ are untagged. The equation $v_1 = u_1^o$ forces $u_1 \in V \sm W$ (else $u_1 \in W$ would give $u_1^o \in W^o$, type-disjoint from $V \sm W$), so $u_1^o = u_1$ and $u_1 = v_1$. Symmetrically $u_2 = v_2$. Hence $(v_1, v_2) = (u_1, u_2) \in L$.

\smallskip\noindent\emph{Conclusion of clause (d).} Combining:
\[
    (v_1, v_2) \in L_{\doit(W)}
        \;\Leftrightarrow\; (v_1, v_2) \in L
        \;\Leftrightarrow\; (v_1, v_2) \in L_{\swig(W)}
        \;\Leftrightarrow\; v_1 \neq v_2 \wedge (v_1, v_2) \in L_{\swig(W)}
        \;\Leftrightarrow\; (\phi(v_1), \phi(v_2)) \in L_{\lp G_{\swig(W)} \rp^{\sm W^o}},
\]
where the third equivalence uses irreflexivity of $L$ (def \ref{def-cdmg}: $(v_1, v_2) \in L \Rightarrow v_1 \neq v_2$) to silently drop or add the conjunct $v_1 \neq v_2$ alongside any $L_{\swig(W)}$-membership statement. \checkmark

\medskip\noindent\emph{Injectivity of the carrier bijection on the source CDMG's node set.}
The canonical bijection $\phi : J_{\doit(W)} \cup V_{\doit(W)} = J \cup V \to J \cup W^i \cup (V \sm W)$ of the statement, defined by $\phi(v) := v$ for $v \in J \cup (V \sm W)$ and $\phi(w) := w^i$ for $w \in W$, sends each element of the source CDMG's joint input-output node set to its image in the target CDMG's carrier $J_{\lp G_{\swig(W)} \rp^{\sm W^o}} \cup V_{\lp G_{\swig(W)} \rp^{\sm W^o}} = J \cup W^i \cup (V \sm W)$ via the three-branch case split spelled out in the statement block. We verify that this bijection is moreover \emph{injective} on the source side -- i.e., distinct elements of $J_{\doit(W)} \cup V_{\doit(W)} = J \cup V$ are sent to distinct elements of $J \cup W^i \cup (V \sm W)$ -- by a case analysis on the three-cell partition of the source. In contrast to claim_3_7's twin (where the analogous injectivity check for two iterated splittings consumes the disjointness side condition $W_1 \cap W_2 = \emptyset$ to rule out two would-be cross-cell collisions), this row's injectivity check needs no extra side condition beyond the lemma's standing hypothesis $W \ins V$: pairwise disjointness of the three source cells follows from def \ref{def-cdmg}'s axiom $J \cap V = \emptyset$ and set algebra alone, and across-cell collisions are then ruled out by def \ref{def:G_node-splitting_intervention}'s tagged-copy disjoint-union construction.

Combining the carrier identifications $J_{\doit(W)} = J \cup W$ and $V_{\doit(W)} = V \sm W$ (items i.\ and ii.\ of def \ref{def:G_hard_intervention}, recalled at the head of the statement block), the source's joint input-output node set decomposes as the disjoint union of three cells
\[
    J_{\doit(W)} \cup V_{\doit(W)} \;=\; J \cup V \;=\; J \;\dcup\; (V \sm W) \;\dcup\; W,
\]
where:
\begin{itemize}
    \item cell (1) $J$ is the input-node set of $G$ (and of $G_{\doit(W)}$ minus $W$);
    \item cell (2) $V \sm W$ is the part of $V$ not affected by the hard intervention (which becomes the output-node set $V_{\doit(W)}$);
    \item cell (3) $W$ is the hard-intervened set of output nodes (which becomes part of the input-node set $J_{\doit(W)}$ as the $W$-summand of $J_{\doit(W)} = J \cup W$).
\end{itemize}
Pairwise disjointness of the three cells follows from the standing CDMG axioms alone: cells (1) and (2) are disjoint by the $J \cap V = \emptyset$ axiom of def \ref{def-cdmg} together with $V \sm W \ins V$; cells (1) and (3) are disjoint because $W \ins V$ (the lemma's side condition) and $J \cap V = \emptyset$ together force $J \cap W = \emptyset$; cells (2) and (3) are disjoint by set algebra ($(V \sm W) \cap W = \emptyset$). \emph{No disjointness side condition is consumed in this partition step} -- the only side condition of the lemma, $W \ins V$, is used only to identify $W$ as a subset of $V$ for the $J$-vs-$W$ disjointness.

\emph{Within-cell injectivity.}
On each of the three cells the bijection $\phi$ acts by a single composed-constructor chain that is injective in the underlying $j$/$v$/$w$-label:
\begin{itemize}
    \item ``$j \mapsto j$'' on cell (1): the identity on $J$, trivially injective.
    \item ``$v \mapsto v$'' on cell (2): the identity on $V \sm W$, trivially injective.
    \item ``$w \mapsto w^i$'' on cell (3): the tagged-copy injection $W \to W^i$ of def \ref{def:G_node-splitting_intervention}, item i.\ (the tagged-input copies $W^i := \lC w^i \st w \in W \rC$ are pairwise distinct as $w$ ranges over $W$, by the type-level disjoint-union construction of the SWIG carrier; in particular $w \mapsto w^i$ is a bijection $W \to W^i$ as recalled at the bijectivity-on-pieces remark of the statement block).
\end{itemize}

\emph{Across-cell injectivity.}
It remains to show that no collision occurs between two elements drawn from two different cells of the three-cell partition. There are $\binom{3}{2} = 3$ structural cross-cell collision patterns to rule out; we walk through each.
\begin{itemize}
    \item \emph{Cell (1) vs.\ cell (2).} Images $\phi(j) = j \in J$ (for some $j \in J$) and $\phi(v) = v \in V \sm W \ins V$ (for some $v \in V \sm W$). An equality $j = v$ inside the codomain $J \cup W^i \cup (V \sm W)$ would force $j = v \in J \cap V = \emptyset$ by the $J \cap V = \emptyset$ axiom of def \ref{def-cdmg} -- contradiction. No collision.
    \item \emph{Cell (1) vs.\ cell (3).} Images $\phi(j) = j \in J$ (for some $j \in J$) and $\phi(w) = w^i \in W^i$ (for some $w \in W$). By def \ref{def:G_node-splitting_intervention}'s tagged-copy construction (items i.\ and ii.\ of def \ref{def:G_node-splitting_intervention} realise the SWIG carrier $J_{\swig(W)} \cup V_{\swig(W)} = (J \dcup W^i) \cup ((V \sm W) \dcup W^o)$ as a type-level disjoint union in which $W^i$ is disjoint from the untagged image of $J \cup (V \sm W)$, with the unsplit-injection shorthand ``$v^i := v$ for $v \in J \cup (V \sm W)$'' fixing the untagged piece), the image-side pieces $J$ and $W^i$ inside the codomain $J \cup W^i \cup (V \sm W)$ are disjoint. No collision.
    \item \emph{Cell (2) vs.\ cell (3).} Images $\phi(v) = v \in V \sm W$ (for some $v \in V \sm W$) and $\phi(w) = w^i \in W^i$ (for some $w \in W$). Same type-disjointness argument as cells (1) vs.\ (3): the untagged image $V \sm W$ and the tagged-input piece $W^i$ are disjoint inside the codomain $J \cup W^i \cup (V \sm W)$ by def \ref{def:G_node-splitting_intervention}'s tagged-copy disjoint-union construction (the SWIG carrier identifies $W^i$ as type-disjoint from every element of $V$, hence in particular from $V \sm W$). No collision.
\end{itemize}
All three cross-cell collision patterns have been ruled out, so the canonical carrier bijection $\phi$ is injective on $J_{\doit(W)} \cup V_{\doit(W)} = J \cup V$. \checkmark

\medskip
Combining the four clauses (a)--(d) with the carrier-bijection injectivity check just established, $\phi$ is a bijection on nodes that respects the input/output partition (clauses (a) and (b)), whose forward and inverse direction biconditionals characterise the directed-edge sets (clause (c)) and the bidirected-edge sets (clause (d)) of the two CDMGs, and which is moreover injective on the source CDMG's joint input-output node set $J_{\doit(W)} \cup V_{\doit(W)} = J \cup V$. Hence $\phi$ is a CDMG-isomorphism (def \ref{def-cdmg}), establishing $G_{\doit(W)} \cong \lp G_{\swig(W)} \rp^{\sm W^o}$.
\end{proof}

\end{document}
